% !PRINTED 426
\weberpage{426}

%<<CONT
gemacht werden k\"onnen, deren Grundz\"uge passend hier eingeschoben
werden, wo sich die Analogie mit der Theorie der Einheitswurzeln
deutlich zeigt.

Es sei also jetzt $n$ eine \gesperrt{Primzahl} und
\weq{1}{f(x) = a_0 x^m + a_1 x^{m-1} + \cdots + a_{m-1} x + a_m}
eine ganze Function von $x$, deren Co\"efficienten $a$ ganze Zahlen
sind, und $a_0$ nicht durch $n$ theilbar. Setzen wir f\"ur $x$ eine
solche ganze Zahl $\alpha$, dass $f(\alpha)$ durch $n$ theilbar
wird, so sagen wir (nach Analogie der Gleichungen), $\alpha$ sei
eine Wurzel der Congruenz $m^{\text{ten}}$ Grades:
\weq{2}{f(x) \equiv 0 \quad (\mathrm{mod}\ n).}
Wird $\alpha$ um ein Vielfaches von $n$ vermehrt, so bleibt es
Wurzel der Congruenz (2). Solche nach dem Modul $n$ congruente
Wurzeln gelten nicht als verschieden. Unter dieser Voraussetzung
k\"onnen wir den Satz aussprechen:

\gesperrt{1. Eine Congruenz $m^{\text{ten}}$ Grades f\"ur einen
Primzahlmodul kann nicht mehr als $m$ verschiedene Wurzeln haben.}

Der Satz ist richtig f\"ur $m = 1$, denn die Congruenz $a_0 x + a
\equiv 0\ (\mathrm{mod}\ n)$ hat nur eine Wurzel, n\"amlich, wenn
% [sic] The source prints modulus m here although the inverse is being taken modulo n; preserved diplomatically.
$a_0'$ so bestimmt wird, dass $a_0 a_0' \equiv 1\ (\mathrm{mod}\ m)$
ist, $\alpha \equiv -a a_0'\ (\mathrm{mod}\ n)$. Wir nehmen unseren
Satz also jetzt als bewiesen an f\"ur den Grad $m-1$ und leiten
seine Richtigkeit f\"ur den Grad $m$ daraus her. Nach \S.\,4 k\"onnen
wir, wenn zun\"achst $\alpha$ beliebig ist, setzen
\weq{3}{\frac{f(x) - f(\alpha)}{x - \alpha} = f_1(x),}
worin $f_1(x)$ eine ganze Function vom $(m-1)^{\text{ten}}$ Grade
ist, die, wenn $\alpha$ eine ganze Zahl ist, ganzzahlige
Co\"efficienten hat. Wenn also $f(\alpha) \equiv 0$ ist, so folgt
aus (3)
\weq{4}{f(x) \equiv (x - \alpha)\, f_1(x) \quad (\mathrm{mod}\ n).}
Jede Wurzel $\beta$ der Congruenz (2) muss also der Bedingung
$(\beta - \alpha)\, f_1(\beta) \equiv 0\ (\mathrm{mod}\ n)$
gen\"ugen. Also ist entweder $\beta - \alpha$ oder $f_1(\beta)$
durch $n$ theilbar.

Nun giebt es nach Voraussetzung h\"ochstens $m-1$ Werthe $\beta$,
f\"ur die $f_1(\beta)$ durch $n$ theilbar wird; giebt es also noch
eine $m^{\text{te}}$ Wurzel von (4), so muss diese gleich $\alpha$
sein.

Wenn also eine Congruenz von der Form $f(x) \equiv 0\ (\mathrm{mod}\
n)$ mehr Wurzeln hat, als ihr Grad betr\"agt, so schliessen wir,
dass
%CONT>>
